% ROZDZIAŁ 27 -- ujęcie miarowe; definicje geometryczne w osobnym skrypcie.
\chapter[Miary na rozmaitościach]{Miary na rozmaitościach: ujęcie miarowo-teoretyczne}
\label{ch:miary-rozmaitosci-27}

\section{Zakres rozdziału i ustalone obiekty}
W opracowaniu \emph{Geometria i topologia różniczkowa}, w rozdziale
\emph{Tensor metryczny, długość i miara}, skonstruowano już
miarę riemannowską, wyprowadzono jej niezależność od mapy,
wprowadzono metrykę indukowaną i obliczono całki po gładkich
podrozmaitościach. W rozdziale o formach udowodniono twierdzenie
Stokesa. Nie powtarzamy tych dowodów ani przykładów
sfery i paraboloidy. Zamiast tego badamy miarę riemannowską
jako obiekt teorii miary: regularność, zbiory zerowe,
gęstości Radona--Nikodyma, transport miar i produkty.

Przez $(M^n,g)$ rozumiemy w tym rozdziale niepustą, gładką,
Hausdorffowską i drugą przeliczalną rozmaitość bez brzegu,
wymiaru $n\ge1$, wyposażoną w gładką metrykę riemannowską.
Nie wymagamy orientowalności. Symbol $\mathcal B(M)$ oznacza
$\sigma$-algebrę borelowską, a $\mu_g$ oznacza miarę
borelowską określoną lokalnie przez
\begin{equation}\label{eq:g-density-27}
  \mu_g(A\cap U)=\int_{x(A\cap U)}
  \sqrt{\det G(x^{-1}(u))}\,\mathrm d\lambda_n(u),
\end{equation}
gdzie $(U,x)$ jest mapą, $A\in\mathcal B(M)$,
a $G=(g_{ij})$ jest macierzą metryki w tej mapie.
Jest to przypomnienie definicji, nie nowa konstrukcja.
Jej uzupełnienie będziemy oznaczać przez
$(\overline{\mathcal B(M)}^{\,\mu_g},\overline\mu_g)$.

\section{Regularność i $\sigma$-skończoność}
\begin{twierdzenie}[Miara objętościowa jest miarą Radona]
\label{thm:riemann-radon-27}
Niech $(M^n,g)$ spełnia powyższe założenia. Wówczas
$\mu_g$ jest lokalnie skończoną, $\sigma$-skończoną
miarą Radona na $\mathcal B(M)$. W szczególności
$\mu_g(K)<\infty$ dla każdego zwartego $K\subseteq M$.
\end{twierdzenie}
\begin{proof}
Weźmy punkt $p\in M$ i mapę $(U,x)$ wokół $p$.
Wybieramy kulę euklidesową $B$ tak, aby
$\overline B\subset x(U)$, a jej przeciwobraz $K=x^{-1}(\overline B)$
był zwartym otoczeniem punktu $p$. Gęstość
$\rho(u)=\sqrt{\det G(x^{-1}(u))}$ jest dodatnia
i ciągła, zatem na $\overline B$ jest ograniczona
z góry i z dołu przez dodatnie stałe. Z
\eqref{eq:g-density-27} wynika $\mu_g(K)<\infty$.
Dzięki zwartości dowolny zwarty $L\subset M$ pokrywamy
skończenie wieloma takimi otoczeniami, skąd
$\mu_g(L)<\infty$.

Rozmaitość $M$ jest lokalnie zwarta i ma przeliczalną
bazę, więc jest $\sigma$-zwarta: wybieramy
przeliczalnie wiele otoczeń o zwartych domknięciach,
które ją pokrywają. W konsekwencji $\mu_g$ jest
$\sigma$-skończona. Ponadto $M$ jest przestrzenią
metryzowalną i lokalnie zwartą. Stosujemy
udowodnione wcześniej twierdzenie o regularności
lokalnie skończonych miar borelowskich na takich
przestrzeniach (Rozdział~\ref{ch:radon-measures-20}),
aby otrzymać regularność wewnętrzną przez zbiory
zwarte i zewnętrzną przez zbiory otwarte.
\end{proof}

\begin{uwaga}[Miara borelowska a jej uzupełnienie]
Nie należy utożsamiać $\mathcal B(M)$ z jej
uzupełnieniem. Podzbiory borelowskiego zbioru
miary zero mogą nie być borelowskie, ale należą do
$\overline{\mathcal B(M)}^{\,\mu_g}$.
Definicja miary Radona dotyczy w tym rozdziale
miary borelowskiej; do całkowania funkcji
mierzalnych względem jej uzupełnienia używamy
standardowego rozszerzenia całki.
\end{uwaga}

\section{Zbiory o zerowej objętości}
\begin{twierdzenie}[Lokalne kryterium zbiorów zerowych]
\label{thm:null-local-27}
Niech $A\in\mathcal B(M)$, a
$\{(U_j,x_j)\}_{j\ge1}$ będzie przeliczalnym atlasem
pokrywającym $M$. Wówczas
\[
 \mu_g(A)=0\quad\Longleftrightarrow\quad
 \lambda_n\bigl(x_j(A\cap U_j)\bigr)=0
 \quad\text{dla każdego }j\ge1.
\]
To samo kryterium działa dla zbiorów mierzalnych
w uzupełnionej $\sigma$-algebrze, gdy ich obrazy
rozumiemy jako zbiory mierzalne w sensie Lebesgue'a.
\end{twierdzenie}
\begin{proof}
Na każdej mapie gęstość $\rho_j$ jest dodatnia
i ciągła. Dla zbioru borelowskiego $E\subseteq x_j(U_j)$
mamy $\int_E\rho_j\,\mathrm d\lambda_n=0$
wtedy i tylko wtedy, gdy $\lambda_n(E)=0$:
wystarcza wyczerpać $x_j(U_j)$ zbiorami zwartymi,
na których $\rho_j$ ma dodatnie minimum.
Z \eqref{eq:g-density-27} otrzymujemy kryterium
lokalne, a przeliczalność atlasu pozwala
zsumować je na całym $M$. W wersji zupełnej
zbiory zerowe można zastąpić borelowskimi
zbiorami zerowymi, które je zawierają.
\end{proof}

\begin{wniosek}[Obrazy immersji niższego wymiaru]
\label{cor:submanifold-null-27}
Niech $N^k$ będzie drugą przeliczalną gładką
rozmaitością, $0\le k<n$, a $F:N\to M^n$
immersją klasy $C^1$. Jeśli $S=F(N)$ jest
zbiorem borelowskim w $M$, to $\mu_g(S)=0$.
W szczególności teza zachodzi dla każdej osadzonej
podrozmaitości gładkiej wymiaru $k<n$.
\end{wniosek}
\begin{proof}
Twierdzenie o stałej randze pozwala dla każdego
$p\in N$ wybrać otoczenie $V_p\subseteq N$,
na którym $F$ jest osadzeniem, a $F(V_p)$
ma w odpowiedniej mapie celu postać kawałka
$\mathbb R^k\times\{0\}^{n-k}$.
Taki zbiór ma zerową $n$-wymiarową miarę
Lebesgue'a z twierdzenia Tonellego, więc
ma zerową $\mu_g$-miarę w każdej mapie celu.
Ponieważ $N$ ma przeliczalną bazę, z pokrycia
$(V_p)_{p\in N}$ wybieramy podpokrycie
przeliczalne $(V_j)$. Zatem
$S=\bigcup_jF(V_j)$, a przeliczalna
podaddytywność miary zewnętrznej daje
$\mu_g^*(S)=0$. Symbol $\mu_g^*$ oznacza tu miarę
zewnętrzną wyznaczoną przez $\mu_g$, czyli infimum
$\mu_g(B)$ po borelowskich $B\supseteq S$. Założenie borelowskości $S$
pozwala zapisać wniosek jako $\mu_g(S)=0$.
Argument nie zakłada, że \emph{cały} obraz immersji
jest lokalnie osadzoną podrozmaitością w otoczeniu
każdego punktu obrazu: przy samoprzecięciach lub
nagromadzeniach może to być nieprawdą.
\end{proof}

\begin{uwaga}[Dwie różne miary na tym samym zbiorze]
Wniosek nie oznacza, że długość krzywej lub pole
powierzchni znika. Dla osadzonej $k$-wymiarowej
podrozmaitości $S$, miara $\mu_{g|_S}$ może być
dodatnia, choć $\mu_g(S)=0$, gdy $k<n$.
To porównanie prowadzi do miar Hausdorffa;
identyfikację $\mu_{g|_S}$ z odpowiednio
unormowaną $\mathcal H^k$ uzasadnimy
w Rozdziale~\ref{ch:area-coarea-28}.
\end{uwaga}

\section{Miary z gęstością i rozkład Lebesgue'a}
\begin{twierdzenie}[Miary absolutnie ciągłe względem objętości]
\label{thm:rn-manifold-27}
Niech $(M,g)$ będzie jak wyżej i niech
$\nu$ będzie dodatnią, $\sigma$-skończoną
miarą na $\mathcal B(M)$, przy czym
$\nu\ll\mu_g$. Wtedy istnieje nieujemna funkcja
borelowska $h$, jednoznaczna $\mu_g$-prawie
wszędzie, taka że dla każdego $A\in\mathcal B(M)$
\[
 \nu(A)=\int_A h(p)\,\mathrm d\mu_g(p).
\]
Jeżeli dodatkowo $\nu$ jest lokalnie skończona,
to $h\in L^1_{\mathrm{loc}}(M,\mu_g)$,
tj. $\int_K h\,\mathrm d\mu_g<\infty$
dla każdego zwartego $K\subseteq M$.
\end{twierdzenie}
\begin{proof}
Miara $\mu_g$ jest $\sigma$-skończona na mocy
Twierdzenia~\ref{thm:riemann-radon-27}.
Stosujemy twierdzenie Radona--Nikodyma
(Rozdział~\ref{ch:radon-nikodym-18})
do pary $(\nu,\mu_g)$. Otrzymujemy $h$ mierzalną
względem uzupełnienia; można wybrać przedstawiciela
borelowskiego, modyfikując ją na zbiorze
$\mu_g$-miary zero. Jeżeli $\nu$ jest lokalnie
skończona, każdy zwarty $K$ można przykryć
skończenie wieloma otoczeniami o skończonej
$\nu$-mierze, więc $\nu(K)<\infty$.
Wstawienie $A=K$ daje lokalną całkowalność.
\end{proof}

\begin{przyklad}[Gęstość w mapie]
Jeżeli w mapie $x:U\to V\subseteq\mathbb R^n$
miara $\nu$ ma gęstość $q(u)$ względem $\lambda_n$,
to gęstość względem $\mu_g$ wynosi
$h(p)=q(x(p))/\sqrt{\det G(p)}$
na $U$, w sensie równości prawie wszędzie.
Wyrażenie $q$ zmienia się po zmianie mapy;
$h$ jest funkcją skalarną na rozmaitości.
\end{przyklad}

\begin{twierdzenie}[Rozkład na część objętościową i osobliwą]
\label{thm:leb-manifold-27}
Jeżeli $\nu$ jest dodatnią, $\sigma$-skończoną
miarą na $\mathcal B(M)$, to ma jednoznaczny rozkład
$\nu=h\mu_g+\nu_{\mathrm s}$, gdzie
$h\ge0$ jest borelowska, a
$\nu_{\mathrm s}\perp\mu_g$.
\end{twierdzenie}
\begin{proof}
Jest to zastosowanie twierdzenia o rozkładzie
Lebesgue'a (Rozdział~\ref{ch:ac-singular-17})
i twierdzenia Radona--Nikodyma do części
absolutnie ciągłej. Jednoznaczność obu składników
wynika z tych twierdzeń.
\end{proof}

\section{Przenoszenie objętości przez dyfeomorfizm}
\begin{definicja}[Jakobian objętościowy]
\label{def:volume-jacobian-27}
Niech $(M^n,g)$ i $(N^n,h)$ będą rozmaitościami
riemannowskimi, a $F:M\to N$ dyfeomorfizmem $C^1$.
Dla $p\in M$ wybieramy ortonormalne bazy
w $T_pM$ i $T_{F(p)}N$, a przez $A_p$ oznaczamy
macierz $\mathrm dF_p$ w tych bazach. Definiujemy
$J_{g,h}F(p)=|\det A_p|$. Wartość nie zależy
od wyboru baz. W lokalnych współrzędnych $u$ i $v$
wyraża się wzorem
\[
 J_{g,h}F(p)=|\det D(v\circ F\circ u^{-1})(u(p))|
 \frac{\sqrt{\det H(v(F(p)))}}{\sqrt{\det G(u(p))}}.
\]
\end{definicja}

\begin{twierdzenie}[Zamiana zmiennych na rozmaitościach]
\label{thm:manifold-change-27}
Dla powyższego dyfeomorfizmu i każdej funkcji
borelowskiej $f:N\to[0,+\infty]$ zachodzi
\begin{equation}\label{eq:manifold-change-27}
 \int_N f(q)\,\mathrm d\mu_h(q)=
 \int_M f(F(p))J_{g,h}F(p)\,\mathrm d\mu_g(p).
\end{equation}
Ta sama równość zachodzi dla funkcji ze znakiem,
gdy $f\in L^1(N,\mu_h)$.
W szczególności $(F_*\mu_g)\ll\mu_h$ i
\[
 \frac{\mathrm d(F_*\mu_g)}{\mathrm d\mu_h}(q)
 =\frac{1}{J_{g,h}F(F^{-1}(q))}
 \quad\text{dla }\mu_h\text{-prawie wszystkich }q.
\]
\end{twierdzenie}
\begin{proof}
Najpierw załóżmy, że $f$ ma nośnik zawarty
w mapie celu, której przeciwobraz przez $F$
mieści się w mapie źródła. Z definicji
$\mu_g,\mu_h$ oraz twierdzenia o zamianie
zmiennych w $\mathbb R^n$ otrzymujemy
\eqref{eq:manifold-change-27}, po skróceniu
wyznaczników macierzy metryk.

W ogólności wybieramy przeliczalne pokrycie
mapami $V_j$ przestrzeni $N$ i rozłączny podział
borelowski $B_1=V_1$, $B_j=V_j\setminus
\bigcup_{i<j}V_i$. Każdy $B_j$ rozbijamy dalej
na przeliczalną rodzinę części, na których
$F^{-1}$ ma wspólną mapę źródła. Dla
$f\chi_{B_j}$ stosujemy argument lokalny,
a następnie sumujemy po rozłącznych częściach,
używając twierdzenia Tonellego.

Aby obliczyć miarę obrazową, przyjmujemy
$f=\chi_A$ w równaniu zastosowanym do $F^{-1}$.
Ponieważ
$J_{h,g}(F^{-1})(q)=[J_{g,h}F(F^{-1}q)]^{-1}$,
dostajemy podaną pochodną Radona--Nikodyma.
Dla funkcji całkowalnych stosujemy równość
oddzielnie do części dodatniej i ujemnej.
\end{proof}

\begin{przyklad}[Konforemna zmiana metryki]
Niech $\varphi\in C^\infty(M;\mathbb R)$ i
$g'=e^{2\varphi}g$. Macierze spełniają
$G'=e^{2\varphi}G$, więc
$\det G'=e^{2n\varphi}\det G$. W konsekwencji
\[
 \mathrm d\mu_{g'}(p)=e^{n\varphi(p)}\,\mathrm d\mu_g(p).
\]
Nie oznacza to, że krzywizna obu metryk jest
taka sama: wzór dotyczy tylko miary objętościowej.
\end{przyklad}

\section{Miara na produkcie rozmaitości}
\begin{twierdzenie}[Miara produktu riemannowskiego]
\label{thm:product-riemann-27}
Niech $(M^n,g)$ i $(N^m,h)$ będą rozmaitościami
jak wyżej. Na $M\times N$ rozważmy metrykę
$g\oplus h$. Wtedy na $\mathcal B(M\times N)$
zachodzi $\mu_{g\oplus h}=\mu_g\otimes\mu_h$.
Dla nieujemnej funkcji borelowskiej $f$ mamy
\[
 \int_{M\times N} f(p,q)\,\mathrm d\mu_{g\oplus h}(p,q)
 =\int_M\!\left(\int_N f(p,q)\,\mathrm d\mu_h(q)\right)
   \mathrm d\mu_g(p).
\]
\end{twierdzenie}
\begin{proof}
W mapach produktowych macierz metryki jest
blokowo diagonalna $\operatorname{diag}(G,H)$.
Stąd $\sqrt{\det(G\oplus H)}=
\sqrt{\det G}\sqrt{\det H}$. Miara lokalna
jest produktem odpowiednich miar lokalnych
na mapach. Równość na całym produkcie wynika
z przeliczalnego rozbicia na mapy i jednoznaczności
miary produktowej dla miar $\sigma$-skończonych.
Ostatnia równość jest twierdzeniem Tonellego.
\end{proof}

\section{Od gładkich podrozmaitości do zbiorów prostowalnych}
W geometrii rozważaliśmy miarę indukowaną
na \emph{gładkich} podrozmaitościach. Obraz
odwzorowania Lipschitzowskiego nie musi być
gładką podrozmaitością, lecz nadal może
posiadać sensowną miarę $k$-wymiarową.
W następnym rozdziale pokażemy, jak
$\mathcal H^k$ mierzy takie obrazy i dlaczego
przy odwzorowaniach nieiniektywnych konieczny
jest czynnik krotności. Wprowadzone tam wzory
są ściśle miarowe: nie potrzebują globalnej
orientacji ani globalnych parametryzacji.